\subsection{Depth First Search (DFS)}
\textbf{Ý tưởng:} DFS cũng là một thuật toán thuộc nhóm \textbf{Uninformed-Search} với cơ chế mở rộng nút sâu nhất chưa được mở rộng; có thể cài bằng đệ quy hoặc ngăn xếp.

\subsubsection{DFS by Recursion}

\textbf{Cấu trúc dữ liệu:} đệ quy trên cây tìm kiếm, ánh xạ \texttt{visited[node] = parent}.

\textbf{Cơ chế thực hiện:}
\begin{itemize}
    \item DFS đệ quy mở rộng nút sâu nhất có thể trước khi quay lui (backtrack), đảm bảo đi đến tận cùng của một nhánh trước khi khám phá các nhánh khác.
    \item Triển khai sử dụng \textbf{gọi hàm đệ quy} thay cho cấu trúc dữ liệu ngăn xếp.
    \item Khi một nút được mở r, thuật toán đánh dấu nút đó là đã \texttt{visited}, sau đó gọi đệ quy để tiếp tục duyệt các láng giềng chưa thăm của nó.
    \item Việc quay lui xảy ra khi không còn láng giềng nào chưa được thăm trong nhánh hiện tại.
    \item Kiểm tra mục tiêu thường được áp dụng ngay khi nút được sinh ra hoặc ngay trước khi gọi đệ quy.
    \item Cơ chế gọi đệ quy mô phỏng chính xác quá trình \textbf{LIFO (Last-In, First-Out)} – tương tự như sử dụng ngăn xếp.
\end{itemize}

\textbf{Ghi chú triển khai:}
\begin{itemize}
    \item DFS đệ quy đơn giản, dễ đọc nhưng dễ gặp \textbf{lỗi tràn ngăn xếp (stack overflow)} vì phải gọi đệ quy nhiều lần với đồ thị sâu.
\end{itemize}

\subsubsection{DFS by Stack}

\textbf{Cấu trúc dữ liệu:} \texttt{explicit stack}, ánh xạ \texttt{visited[node] = parent}.

\textbf{Cơ chế thực hiện:}
\begin{itemize}
    \item Phiên bản này thay thế cơ chế đệ quy bằng một \textbf{ngăn xếp tường minh} để lưu các nút chờ được mở rộng.
    \item Thuật toán bắt đầu bằng cách push nút gốc vào ngăn xếp.
    \item Ở mỗi bước, pop một nút trên đỉnh ngăn xếp, đánh dấu là \texttt{visited}, sau đó push các láng giềng chưa thăm của nó vào ngăn xếp.
    \item Kiểm tra mục tiêu có thể được thực hiện ngay sau khi pop nút ra khỏi ngăn xếp.
\end{itemize}

\textbf{Ghi chú triển khai:}
\begin{itemize}
    \item Thứ tự push các láng giềng được sắp xếp \textbf{tăng dần} để khi pop, các nút có chỉ số lớn hơn (nhánh phải) được xử lý trước — phù hợp quy tắc “ưu tiên nhánh phải”.
    \item DFS bằng ngăn xếp kiểm soát tốt hơn bộ nhớ và thứ tự duyệt nhưng tiêu tốn bộ nhớ hơn, phù hợp cho mô phỏng và trực quan hóa.
\end{itemize}

\textbf{Đánh giá hiệu suất:} \cite{lecture2025}
\begin{itemize}
    \item \textbf{Tính đầy đủ (Completeness):} Có, nếu không gian trạng thái hữu hạn và có kiểm soát vòng lặp.
    \item \textbf{Tính tối ưu (Optimality):} Không, vì DFS có thể tìm thấy lời giải “bên trái hoặc phải nhất” (tùy theo cách cài đặt) thay vì ngắn nhất.
    \item \textbf{Độ phức tạp thời gian (Time Complexity):} $O(b^m)$, trong đó $b$ là hệ số phân nhánh và $m$ là độ sâu tối đa của không gian trạng thái.
    \item \textbf{Độ phức tạp không gian (Space Complexity):} 
    \begin{itemize}
        \item $O(m)$ đối với DFS đệ quy (chỉ giữ ngăn xếp lời gọi hàm).
        \item $O(bm)$ đối với DFS bằng ngăn xếp.
    \end{itemize}
\end{itemize}
